\chapter{Evaluation Design}
\label{ch:evaluation}

The evaluation is the thesis's graded core, and this chapter defines it before any numbers
are reported. We cover the financial metrics and their regime segmentation
(\cref{sec:finmetrics}), then the three reasoning metrics --- faithfulness
(\cref{sec:faithfulness}), grounding (\cref{sec:grounding}), and sophistication
(\cref{sec:sophistication}) --- explaining in each case \emph{why} the chosen measurement
strategy fits the quantity. \Cref{sec:anchor} treats the judge--human anchoring that
underwrites the validity of the faithfulness metric. \Cref{sec:pmeval} extends the protocol
to the multi-agent path, and \cref{sec:rigour} presents the rigour narrative: the chain of
bugs caught, diagnosed, fixed, and tested. \Cref{fig:evalprotocol} summarises the protocol.

% ---------------------------------------------------------------------
% Figure: evaluation protocol. A logged decision fans out to three
% measurements: faithfulness (blind LLM judge), grounding (deterministic),
% and the human anchor that validates the judge.
% ---------------------------------------------------------------------
\begin{figure}[t]
  \centering
  \begin{tikzpicture}[
      >=Stealth, font=\footnotesize, node distance=12mm,
      rec/.style={draw, rounded corners, align=center, minimum height=12mm,
                  minimum width=30mm, thick, fill=gray!8},
      m/.style={draw, rounded corners, align=center, minimum height=12mm,
                minimum width=34mm, thick},
      judge/.style={m, fill=accent!8, draw=accent!55},
      det/.style={m, fill=accent!8, draw=accent!55},
      hum/.style={m, fill=todored!8, draw=todored!55},
    ]
    \node[rec] (dec) {Logged decision\\\scriptsize evidence + \CoT\\\scriptsize + executed order};

    \node[judge, above right=8mm and 26mm of dec] (faith)
      {\textbf{Faithfulness}\\\scriptsize blind judge predicts\\\scriptsize rational action (temp 0)};
    \node[det, below right=8mm and 26mm of dec] (ground)
      {\textbf{Grounding}\\\scriptsize numeric claims vs\\\scriptsize actual tool outputs ($\pm1\%$)};

    \node[hum, right=22mm of faith] (human)
      {\textbf{Human anchor}\\\scriptsize blind annotation\\\scriptsize Cohen's $\kappa$};

    % what the judge sees vs what the agent did
    \draw[->, thick] (dec.east) -- ++(8mm,0) |- (faith.west)
      node[pos=0.75, above, font=\scriptsize]{evidence only};
    \draw[->, thick] (dec.east) -- ++(8mm,0) |- (ground.west)
      node[pos=0.75, below, font=\scriptsize]{\CoT\ + outputs};
    \draw[->, thick, dashed, todored!70] (faith) -- (human)
      node[midway, above, font=\scriptsize]{calibrate};

    \node[align=center, below=7mm of ground, text width=80mm, font=\scriptsize\itshape,
          text=gray]
      {compare judge prediction to the agent's \emph{executed} action $\Rightarrow$
       faithful / constrained / unfaithful, then segment every metric by market regime};
  \end{tikzpicture}
  \caption[Reasoning-evaluation protocol]{The reasoning-evaluation protocol. Each logged
  decision is measured along two automatic axes. \emph{Faithfulness} asks whether the
  executed trade matches the action a rational, policy-aware judge would take given
  \emph{only} the same objective evidence (regime, indicators, portfolio) --- deliberately
  withholding the agent's own rationale and action so the judge cannot simply echo them.
  \emph{Grounding} is deterministic: it extracts numeric claims from the chain-of-thought
  and checks each against the actual tool outputs within a $1\%$ tolerance. The judge is in
  turn anchored to blind human annotation (Cohen's $\kappa$), making the human agreement the
  validity cornerstone of the whole protocol. Every metric is finally segmented by regime.}
  \label{fig:evalprotocol}
\end{figure}


\section{Financial performance}
\label{sec:finmetrics}

Financial performance is reported against the baselines of \cref{sec:baselines} using a
standard risk-adjusted metric set computed in \file{eval/financial.py}: annualised return,
the Sharpe~\cite{sharpe1994} and Sortino~\cite{sortino1994} ratios, the Calmar ratio,
maximum drawdown, and hit rate, alongside turnover and cost drag. The metrics are pure
functions of the \nav\ series, which makes them trivially testable and cacheable. Each point
estimate is accompanied by the bootstrap confidence interval of \cref{sec:bootstrap}, and
the agent is compared on the \emph{same} window, the \emph{same} costs, and the \emph{same}
shift discipline as the baselines.

Crucially, financial performance is reported \emph{by regime}. A headline Sharpe that
aggregates a bull run with a high-volatility episode hides exactly the behaviour the thesis
is about. The regime labels of \cref{sec:regime} partition the decision sequence, and every
metric is recomputed within each partition. The expected --- and most interesting ---
finding is that performance and reasoning quality both degrade in adversarial regimes, in
line with ATLAS~\cite{atlas2025} and LiveTradeBench~\cite{livetradebench2025}.

\section{Faithfulness: the knowledge--action gap}
\label{sec:faithfulness}

\subsection{What it measures, and why a judge}

Faithfulness measures the knowledge--action gap of KellyBench~\cite{kellybench2026}: does the
agent \emph{do} what its reasoning implies it should do? The naive implementation --- parse
an intended action from the rationale and compare it to the executed action --- measures
almost nothing, because both fields are produced by the same LLM call and are trivially
self-consistent. Keyword extraction from the rationale is also fragile on nuanced prose
(``I will not sell yet'' misfires). The faithfulness signal we want is not internal
self-consistency but consistency between the agent's behaviour and what a \emph{rational
observer} would do on the same evidence.

We therefore use an independent LLM judge (temperature $0$ for determinism) that sees
\emph{only} the objective evidence the agent saw --- ticker, date, regime, technical
indicators, and the portfolio snapshot --- and predicts the rational action \emph{blind}: it
is shown neither the agent's rationale nor its executed action, so it cannot echo them. We
then compare the judge's prediction to the agent's actual executed action. This is the
reason faithfulness uses a judge while grounding does not (\cref{sec:grounding}):
faithfulness is an inherently judgemental counterfactual --- ``what \emph{should} a rational
agent have done here?'' --- whereas grounding is a factual check that has a deterministic
ground truth.

\subsection{Three verdicts}

Each scored decision receives one of three verdicts:
\begin{description}
  \item[Faithful] the agent's action matches what the policy-aware judge would do --- whether
  that is to trade or, deliberately, to hold.
  \item[Constrained] the agent could not have acted: a genuine \emph{capacity} constraint
  (e.g.\ cash below a minimal-trade threshold on the buy side, or a zero position on the
  sell side). These are excluded from the strict faithfulness denominator because they are
  not reasoning failures.
  \item[Unfaithful] a real gap: the judge predicts an action, the agent does not take it, and
  no capacity reason explains the difference.
\end{description}
The standard faithfulness rate counts faithful over all scoreable decisions; a
\emph{strict} variant additionally removes constrained decisions from the denominator. We
report both for comparability.

\subsection{Policy-awareness: the v1$\rightarrow$v2 redesign}
\label{sec:v2judge}

The first judge (v1) assumed a rational agent always deploys to the hard cap. This
systematically mislabelled deliberate holds at the $10$--$15\%$ target as unfaithful,
because the judge ``saw room to buy'' that the agent's own policy forbade using. The
correction (v2) makes the judge \emph{policy-aware}: it is told the base target ($10\%$),
the high-conviction target ($15\%$), and the hard cap ($20\%$), and that once a position is
at or above target, holding is the rational default unless a sell condition exists.

This redesign was \emph{pre-registered}: the policy specification and the verdict semantics
were locked, and validated by the author, \emph{before} the v2 numbers were computed. Locking
the metric definition before seeing its result is a direct response to KellyBench's warnings
about post-hoc metric tuning, and it is what lets us report the v1/v2 comparison as a
genuine methodological finding rather than a fit to a desired number. Both judges are kept in
the codebase behind a flag, and the thesis reports them side by side: v1 for comparison and
reproducibility, v2 as the canonical metric.

A structural consequence of policy-awareness is that the v2 judge resolves capacity cases
itself (it predicts ``hold'' for an at-cap or cash-starved position), so the explicit
constrained-override fires far less often --- the prompt does the work the override used to
do. The redesign also \emph{exposed} a finding the v1 judge had hidden: a class of
non-trimming decisions that v1 silently absorbed into ``constrained'' but v2 correctly marks
unfaithful. This is discussed quantitatively in \cref{ch:results}.

\section{Grounding: are the cited numbers real?}
\label{sec:grounding}

Grounding asks whether the numbers the agent cites in its reasoning are the actual outputs of
the tools it called, or hallucinations. Because this question has a deterministic ground
truth --- the tool outputs are recorded in the decision record --- it is computed
deterministically, with no LLM in the loop. A regex extracts \texttt{(label, value)} numeric
claims from the rationale; each claim is matched against the recorded indicator dictionary,
the recent OHLCV bars, or the portfolio snapshot within a $1\%$ relative tolerance. Grounding
is the fraction of claims that match. Decisions with no numeric claims score as ``no signal''
and are not penalised --- an unquantified rationale is not a hallucination.

The contrast with faithfulness is deliberate and is the methodological point of having both:
grounding is cheap, deterministic, and re-runnable for free on any decision file, so it
anchors the evaluation in a measurement that involves no model judgement at all; faithfulness
necessarily involves judgement and is therefore the metric that must be anchored to humans
(\cref{sec:anchor}). Reporting both protects against the failure mode where an agent cites
perfectly real numbers (high grounding) but acts against them (low faithfulness), which is
exactly the knowledge--action gap.

\section{Sophistication}
\label{sec:sophistication}

The third reasoning axis, planned in the manner of KellyBench's multi-point sophistication
rubric, scores a sampled subset of decisions on a small, explicit rubric --- risk
management, acknowledgement of uncertainty, regime adaptation, and coherence --- via an
LLM judge with structured output, sampled per regime to bound cost. Sophistication is the
most subjective of the three metrics and is reported as a planned, supporting measure rather
than a headline; its rubric maps directly onto the thesis claims about regime-adaptive
reasoning. Its results, where computed, are reported in \cref{ch:results} with the same
regime segmentation and the same candour about sampling.

\section{Anchoring the judge to humans}
\label{sec:anchor}

A metric computed by an LLM judge is only as trustworthy as its agreement with human
judgement; this anchoring is the validity cornerstone of the reasoning evaluation. We draw a
stratified sample of decisions, have them annotated blind against the same policy the judge
uses, and report the judge--human action agreement and Cohen's $\kappa$~\cite{cohen1960}.

Two points of rigour govern how the agreement is reported. First, the sample is
\emph{stratified} (deliberately over-sampling rare and diagnostic cases such as the
non-trimming class), not a representative random draw; we therefore report agreement
\emph{per stratum} and refrain from presenting a single global $\kappa$ as if the sample were
representative. If a population-level estimate is needed, a separate random sample must be
drawn. Second, a blind protocol is used for the medium-window reconciliation: annotations are
filled against sealed fields, and the key is opened only afterward, so the human labels
cannot be retrofitted to the judge. The per-stratum agreement, the global figures with their
representativeness caveat, and the blind-reconciliation result are all reported in
\cref{sec:resultsanchor}.

\section{Extending the protocol to the multi-agent path}
\label{sec:pmeval}

The PM path needs two adaptations. First, the regime label of a portfolio-level decision is
aggregated across its tickers by worst case (\textsc{high\_vol} $>$ \textsc{bear} $>$
\textsc{range} $>$ \textsc{bull}), so a portfolio decision inherits the most adverse regime
it faces. Second, faithfulness is redefined at the allocation level: for each
ticker after the first (baseline) decision, the PM's target-weight \emph{direction}
(increase/hold/decrease) is checked against the weighted analyst consensus, with the same
capacity-constraint exemptions (no cash to rebalance, already concentrated, cannot short from
zero).

Grounding also splits in two. The PM's final rationale is concise and often carries no
numeric claims, so the standard grounding metric can legitimately return ``no claims'' for
it --- which is not a failure. To audit the \emph{analysts'} evidence, a separate
\emph{evidence-grounding} metric checks each analyst's technical and risk claims against
indicators, bars, portfolio, and regime, and each news claim against the causal news tool
outputs by source and publication date. Finally, an artifact-level \emph{MCP time-machine
audit} sweeps the entire decision file for any tool output, order, or fill timestamped after
\tnow, and reports the per-source counts so that cache-versus-API provenance stays visible.
This audit is the empirical counterpart of \cref{inv:nolookahead}: it demonstrates, on the
completed artifact, that no future timestamp ever reached the agent.

\section{The rigour narrative: bugs caught, diagnosed, fixed, tested}
\label{sec:rigour}

This section is an argument, not a confession. A thesis whose central claim is methodological
rigour is strengthened, not weakened, by showing that its instruments \emph{caught} real
defects and that each defect was closed with a regression test. Each item below follows the
same shape: symptom $\rightarrow$ diagnosis $\rightarrow$ fix $\rightarrow$ test.

\begin{description}
  \item[Price-corruption masquerade.] An early run collapsed not from look-ahead but from a
  single-ticker download whose column structure yielded a non-scalar where a scalar close was
  expected; a permissive exception handler accepted the malformed bars silently. The
  anti-look-ahead tests checked \emph{when} a bar was dated, never \emph{what} it contained.
  \emph{Fix}: a fail-loud data-sanity gate (\file{bar\_validation.py}) plus a price-continuity
  guard rejecting implausible adjacent-bar jumps, on both write and read paths. \emph{Test}:
  a suite that reproduces the exact failure mode (multi-element series) and a cache-read
  regression guard.
  \item[Engine cap inverted.] A position-cap clamp set the order quantity \emph{up} to the
  cap rather than \emph{down}, inflating small legitimate orders to the ceiling and destroying
  the portfolio. \emph{Fix}: clamp strictly downward (\code{min(quantity, capped)}).
  \emph{Test}: order-sizing tests that fail if the clamp direction inverts.
  \item[Look-ahead in the regime detector (v1$\rightarrow$v2).] The v1 volatility threshold
  used a full-series percentile, so a future spike could re-label past bars. \emph{Fix}: a
  trailing 252-bar percentile, fully causal. \emph{Test}: an assertion that fails if a
  full-series percentile is reintroduced.
  \item[Asymmetric transaction costs.] The agent paid nothing while baselines paid friction,
  invalidating the comparison. \emph{Fix}: $10$\,bps symmetric costs on both sides plus
  reported cost drag. \emph{Test}: a per-run assertion of cost symmetry.
  \item[News hallucination.] With no news available, the news analyst fabricated sentiment
  from technical indicators. \emph{Fix}: restrict the news prompt to news outputs only and
  emit a zero-weight ``unavailable'' marker when no items exist. \emph{Test}: a check that
  news evidence may not contain technical-indicator keywords, plus consensus-weight tests.
  \item[Silent PM failures.] A broad exception handler turned PM failures into all-cash
  fallbacks indistinguishable from genuine allocations. \emph{Fix}: tag errored decisions
  explicitly and exclude them from every metric, logging the raw completion that failed.
  \emph{Test}: the metrics filter on the error tag.
  \item[Warm-up too short.] The warm-up window was shorter than the volatility-percentile
  window, so early bars used a compressed threshold. \emph{Fix}: derive the warm-up from the
  detector parameters ($271$ trading days) and centralise the constant. \emph{Test}: an
  anti-regression test that fails if the constant is too small or the window changes without
  realignment.
  \item[Token truncation and transient API errors.] Verbose chain-of-thought JSON truncated
  at a small token cap, poisoning the response cache; and a transient server error aborted the
  run before the first valid bar. \emph{Fix}: raise the token cap and retry genuinely
  transient errors with back-off. \emph{Test}: parse-clean assertions on the decision schema.
\end{description}

Taken together these were not random slips but a systematic hardening of every layer ---
data, execution, evaluation, orchestration --- each closed by a test that would catch its
recurrence. The suite that results (\nTests\ tests at the time of writing) is the standing
evidence for the rigour claim, and the most sensitive of these tests --- the clock invariant
and the metric computations --- are the components the supervisor reviews directly.
